\documentclass[10pt]{beamer}

% --- Packages ---
\usepackage{amsmath, amssymb}
\usepackage[authoryear]{natbib}
\bibliographystyle{abbrvnat}

\usepackage{graphicx}
\usepackage{booktabs}

% (Optional) compile with XeLaTeX/LuaLaTeX for fonts
\usepackage{fontspec}
\setmainfont{Times New Roman}
\usepackage{unicode-math}

\usepackage{tikz}
\usetikzlibrary{arrows.meta, positioning, decorations.pathreplacing, calc, fit, backgrounds}

\setbeamertemplate{footline}[frame number]
\setbeamertemplate{navigation symbols}{}
\setbeamertemplate{caption}{\insertcaption}  % keep figure captions plain

\title{Investment Demand Shock and the Business Cycle}
\subtitle{New Results}
\author{Li Xuanguang}
\date{\today}

\begin{document}

\begin{frame}
	\titlepage
\end{frame}

% ==================================================
% Roadmap
% ==================================================
\begin{frame}{What is new in this version}
	Five new results that dig into the baseline mechanism and its fragility:
	\begin{enumerate}
		\item Where does the high \textbf{Labor input} variance come from?
		\item Are the \textbf{inflation and wage} signs (IST vs.\ ID) stable across parameters?
		\item Why is the estimated \textbf{consumption} impact negative, while theory predicts positive?
		\item How does a full \textbf{Smets--Wouters (2007)} model change the picture?
		\item How does an \textbf{i.i.d.\ (independent)} investment demand shock change the result?
	\end{enumerate}

	\vspace{0.6em}
	\footnotesize
	Notation: $\varepsilon_s$ = IST shock, $\varepsilon_d$ = investment demand (ID) shock,
	$\varepsilon_m$ = MEI shock, $\varepsilon_a$ = technology, $\varepsilon_R$ = monetary,
	$\varepsilon_g$ = government, $\varepsilon_b,\varepsilon_p,\varepsilon_w$ = preference / price- / wage-markup.
\end{frame}

% ==================================================
% 1. Labor input
% ==================================================
\section{Labor input}

\begin{frame}{1. Why model hours are too volatile}
	\begin{itemize}
		\item Baseline matches most moments, but $\mathrm{Var}(\hat N_t)$ is far above the data (see model-fit table).
		\item Decompose the log-linear labor demand/supply block:
		\[
			\hat N_t = \frac{\hat{mc}_t + \alpha \hat{k}^s_t + \hat{\mu}_{a,t} - \theta_w \sigma \hat{c}_t}{\alpha + \theta_w \sigma_\ell}.
		\]
		\item $R^2$ of each component on $\hat N_t$: \textbf{consumption co-movement explains $\approx 82\%$} of hours volatility; all other channels are negligible.
	\end{itemize}

	\vspace{1em}
	\centering
	\footnotesize
	\begin{tabular}{l r r}
		\hline\hline
		Component & $\mathrm{corr}(N_t,x_t)$ & $R^2$ (\%) \\ 
		\hline
		Consumption, $\hat{c}_t$ & +0.907 & 82.29 \\
		Marginal cost, $\hat{mc}_t$ & +0.061 & 0.37 \\ 
		Capital services, $\hat{k}_t^s$ & -0.030 & 0.09 \\ 
		Technology, $\hat{\mu}_{a, t}$ & +0.087 & 0.76 \\ 
		\hline
	\end{tabular}
\end{frame}

\begin{frame}{1. The root cause is volatile consumption}
	\begin{itemize}
		\item Model $\mathrm{Var}(\hat c_t)$ (level) is $39.2$ vs.\ $2.8$ in the data, and this is not because consumption is excluded from the baseline data set.
		\item The \textbf{same pathology exists in Smets--Wouters (2007)}: their level variance of consumption is $33.8$ (computed from their replication package).
		\item Difference is hours: SW (2007) fits $\mathrm{Var}(\hat N_t)$ much better because the \textbf{wage markup shock} $\mu_{w,t}$ loosens the wage--$MRS$ link:
		\[
			\hat{\pi}_t^w = \beta \mathbb{E}_t \hat{\pi}_{t+1}^w - \kappa_w \underbrace{\Big(\hat{w}_t - \hat{mrs}_t + \tfrac{1}{\kappa_w}\mu_{w,t}\Big)}_{\text{wage markup}}.
		\]
		\item Without that shock, the wage is nearly static, so consumption volatility passes straight into hours.
	\end{itemize}

	\vspace{0.2em}
	\centering
	{\footnotesize
	\begin{tabular}{lccc}
		\hline\hline
		Variable & \cite{smetsShocksFrictionsUS2007} & Baseline model & Data \\
		\hline
		Consumption ($\hat{c}_t$) & 33.75 & 39.23 & 2.80 \\
		\quad $\rho_1$    & 0.993 & 0.996 & 0.901 \\
		Hours worked ($\hat{N}_t$)& 8.85  & 32.22 & 6.29 \\
		\quad $\rho_1$    & 0.975 & 0.993 & 0.907 \\
		\hline
	\end{tabular}}
\end{frame}

% ==================================================
% 2. Inflation & wage sign stability
% ==================================================
\section{Sign stability}

\begin{frame}{2. Is the inflation/wage sign stable across parameters?}
	\textbf{Question:} the IRFs show \emph{deflationary IST} and \emph{inflationary ID}. How robust are these signs to the parameters?
	\begin{itemize}
		\item Inflation $=$ PDV of marginal cost: $\displaystyle \hat\pi_t = \kappa_p \sum_{j\ge0}\beta^j \mathbb{E}_t[\widehat{mc}_{t+j}]$.
		\item Two marginal-cost channels: wage and capital rental rate.
	\end{itemize}
	\vspace{0.4em}
	\begin{columns}[T]
		\column{0.5\textwidth}
		\centering \textbf{IST shock $\varepsilon_s$}
		\begin{itemize}\small
			\item Inflation sign \textbf{not robust}: flips to inflationary in many cases.
			\item But the \textbf{wage sign is stable} (negative) across the sweep.
		\end{itemize}
		\column{0.5\textwidth}
		\centering \textbf{ID shock $\varepsilon_d$}
		\begin{itemize}\small
			\item Inflation \textbf{robustly positive} across all 500 posterior draws and every parameter sweep.
			\item Wage \textbf{robustly positive} as well.
		\end{itemize}
	\end{columns}
\end{frame}

\begin{frame}{2. IST: inflation sign flips, wage sign does not}
	\centering
	\includegraphics[width=\textwidth]{figures/ist_sign_robustness.pdf}
	\\[0.3em]
	{\footnotesize Each panel varies one parameter over a wide grid. The impact \emph{inflation} sign is fragile; the \emph{wage} sign stays negative.}
\end{frame}

\begin{frame}{2. The role of $\psi_u$: why the IST inflation sign flips}
	\small
	$\psi_u$ is the curvature of the utilization-cost function. The utilization FOC plus the firm's cost minimization give
	\[
		\hat r^k_t = \frac{\psi_u}{1-\psi_u}\,\hat u_t,
		\qquad
		\hat r^k_t = \psi_u\big(\hat w_t + \hat N_t - \hat k_{t-1}\big),
	\]
	so $\hat r^k_t$ moves only a \textbf{fraction $\psi_u$} of $(\hat w_t+\hat N_t)$.
	\begin{itemize}
		\item \textbf{Low $\psi_u$ (baseline $\approx 0.1$):} rental channel muted; inflation is driven by the wage, which falls under IST $\Rightarrow$ \textbf{deflation}.
		\item \textbf{High $\psi_u$:} larger $a''(1)$ $\Rightarrow$ steeper adjustment cost $\Rightarrow$ capital supply can't adjust; rising labor pushes $\hat r^k_t$ up so the rental channel dominates $\Rightarrow$ \textbf{inflation}.
	\end{itemize}
	\vspace{0.2em}
	\centering
	\includegraphics[width=0.92\textwidth]{figures/mc_decomposition_rk_psi_u.pdf}
\end{frame}

\begin{frame}{2. ID: inflation and wage are robustly positive}
	\centering
	\includegraphics[width=0.5\textwidth]{figures/inv_demand_mc_decomposition.pdf}
	\includegraphics[width=0.9\textwidth]{figures/inv_demand_sign_robustness.pdf}

\end{frame}

% ==================================================
% 3. Consumption impact
% ==================================================
\section{Consumption}

\begin{frame}{3. Why estimated consumption response is negative}
	\begin{itemize}
		\item In the estimated baseline, the impact consumption response is negative, and \textbf{no single} parameter change reverses it. A positive response needs a \textbf{combination}.
		\item Aligning with \citeauthor{nireiOriginsPropagationAnimal}'s calibration and varying $\kappa_p$, $\psi_\pi$, $\theta_w$ jointly: consumption turns positive only when \textbf{all three are low}.
	\end{itemize}
	\vspace{0.3em}
	\centering
	\includegraphics[width=0.78\textwidth]{figures/inv_demand_c1_kappa_psipi_heatmap.pdf}
\end{frame}

\begin{frame}{3. The mechanism: kill the short-run substitution effect}
	\textbf{Intuition:} in the long run an ID shock just accumulates capital $\Rightarrow$ households are wealthier $\Rightarrow$ they consume more today. A positive impact response requires \emph{shutting down} the short-run intertemporal substitution channel.
	\vspace{0.5em}
	\begin{itemize}
		\item Low $\theta_w$ $\Rightarrow$ wage barely moves $\Rightarrow$ marginal cost barely moves.
		\item Low $\kappa_p$ (flat Phillips curve) $\Rightarrow$ inflation barely moves.
		\item Low $\psi_\pi$ (inactive Taylor rule) $\Rightarrow$ nominal rate barely moves.
	\end{itemize}
	\vspace{0.5em}
	\centering
	Stable real interest rate $\;\Longrightarrow\;$ no shift of $c_t$ to the future $\;\Longrightarrow\;$ \textbf{positive $\hat c_1$}.
	\\[0.5em]
\end{frame}

% ==================================================
% 4. Smets-Wouters model
% ==================================================
\section{SW (2007) model}

\begin{frame}{4. Smets--Wouters (2007) model: what changes}
	Expand the baseline with: habit formation, monopolistic labor market (Erceg-type wage Phillips curve), and preference / price-markup / wage-markup shocks; add consumption and wage growth as observables.
	\begin{itemize}
		\item \textbf{Better hours fit:} $\mathrm{Var}(\hat N_t)$ drops from $32.2$ to $9.8$ (data $6.3$) --- the wage markup shock does the work.
		\item \textbf{ID shock becomes marginal:} only $4.4\%$ of output and $13.6\%$ of investment FEVD; near-zero for everything else.
		\item Estimated $\kappa_p\approx0.03$, $\kappa_w\approx0.02$ (very flat) --- partly propped up by the markup shocks.
	\end{itemize}
	\vspace{0.2em}
	\centering
	{\tiny% Auto-generated by test_variance_decomp_nocap.m
% file_name: sw_type_no_rel

\begin{table}[!htbp]
\centering
\caption{Unconditional FEVD (\%): Smets--Wouters-style model.}
\label{tab:vd:sw-type-no-rel}
\begin{tabular}{l r r r r r r r r}
\hline
Observable & $\varepsilon_g$ & $\varepsilon_{s}$ & $\varepsilon_R$ & $\varepsilon_a$ & $\varepsilon_b$ & $\varepsilon_p$ & $\varepsilon_w$ & $\varepsilon_{d}$ \\ 
\hline
Output & 6.15 & 10.31 & 7.86 & 43.58 & 17.21 & 5.00 & 5.52 & 4.37 \\ 
Inflation & 0.05 & 1.78 & 0.82 & 18.85 & 0.40 & 53.05 & 25.04 & 0.00 \\ 
Interest rate & 0.19 & 7.16 & 30.37 & 26.67 & 4.51 & 10.39 & 20.72 & 0.00 \\ 
Hours & 2.00 & 20.43 & 7.10 & 34.85 & 7.86 & 1.10 & 26.35 & 0.30 \\ 
Investment & 0.25 & 54.98 & 2.43 & 18.59 & 3.19 & 0.62 & 6.33 & 13.62 \\ 
Wage & 0.00 & 0.12 & 0.14 & 18.25 & 0.30 & 16.96 & 64.23 & 0.00 \\ 
Consumption & 0.25 & 2.36 & 6.48 & 43.49 & 42.01 & 1.92 & 3.49 & 0.00 \\ 
Gov. spending & 99.90 & 0.01 & 0.01 & 0.05 & 0.02 & 0.01 & 0.01 & 0.00 \\ 
\hline
\end{tabular}
\end{table}
}
\end{frame}

\begin{frame}{4. SW model: IRFs flip relative to the baseline}
	\centering
	\includegraphics[width=0.92\textwidth]{figures/compare_irf_sw_type_no_rel.pdf}
	\\[0.3em]
	\begin{itemize}\footnotesize
		\item[\textbf{(i)}] ID shock now generates \emph{much smaller} fluctuations than IST (smaller estimated s.d.).
		\item[\textbf{(ii)}] IST shock now \emph{inflationary}, with a \emph{positive} wage response (sign flip vs.\ baseline).
		\item[\textbf{(iii)}] Consumption response to ID is much smoother: $\hat c_1$ is only $0.2\%$ of $\hat\imath_1$ (vs.\ $12\%$ in baseline).
	\end{itemize}
\end{frame}

\begin{frame}{4. SW model: consumption sign revisited}
	\begin{columns}[c]
		\column{0.55\textwidth}
		\centering
		\includegraphics[width=\textwidth]{figures/inv_demand_c1_psiy_sigmal_heatmap.pdf}
		\column{0.45\textwidth}
		\footnotesize
		Now sweep $\sigma_\ell$ (inverse Frisch) and $\psi_y$ (output-gap coefficient), because the flat Phillips/wage curves already sit near the \citeauthor{nireiOriginsPropagationAnimal} calibration.
		\vspace{0.5em}

		Same logic: a positive $\hat c_1$ needs a low real rate.
		\begin{itemize}
			\item Small $\psi_y$: directly lowers Taylor-rule output-gap feedback.
			\item Small $\sigma_\ell$: suppresses the real wage $\Rightarrow$ smaller wage gap $\Rightarrow$ smaller output gap.
		\end{itemize}
	\end{columns}
	\vspace{0.3em}
\end{frame}

% ==================================================
% 5. Independent ID shock
% ==================================================
\section{i.i.d.\ ID shock}

\begin{frame}{5. Independent (i.i.d.) investment demand shock}
	Households learn the realized investment demand next period and undo it, so the shock may be \emph{transitory}. Set $\rho_d=0$, keep everything else at baseline.
	\begin{itemize}
		\item ID shock contributes \textbf{$22.95\%$ of output} and \textbf{$15.18\%$ of investment} FEVD.
		\item Down from $\approx40\%$ / $\approx50\%$ in the baseline: zero persistence cuts off the future influence.
	\end{itemize}
	\vspace{0.2em}
	\centering
	{\scriptsize% Auto-generated by test_variance_decomp_nocap.m
% file_name: per_add_mei

\begin{table}[!htbp]
\centering
\caption{Variance decomposition (\%).}
\label{tab:vd:per-add-mei}
\begin{tabular}{l r r r r r}
\hline
Observable & $\varepsilon_g$ & $\varepsilon_{s}$ & $\varepsilon_R$ & $\varepsilon_a$ & $\varepsilon_{d}$ \\ 
\hline
Output & 6.67 & 38.60 & 9.35 & 22.42  & 22.95 \\ 
Inflation & 0.40 & 52.00 & 12.87 & 34.37 & 0.37 \\ 
Interest rate & 0.37 & 74.14 & 3.74 & 21.66 & 0.09 \\ 
Hours & 2.84 & 92.22 & 0.12 & 4.53 & 0.29 \\ 
Investment & 0.25 & 73.06 & 0.20 & 11.31 & 15.18 \\ 
Gov. spending & 99.59 & 0.17 & 0.04 & 0.10 & 0.10 \\ 
\hline
\end{tabular}
\end{table}
}
\end{frame}

% ==================================================
% Wrap-up
% ==================================================
\begin{frame}{Summary of new results}
	\begin{enumerate}
		\item \textbf{Hours:} excess volatility comes almost entirely ($R^2\approx82\%$) from over-volatile consumption; a wage markup shock cures the hours fit.
		\item \textbf{Signs:} ID inflation/wage are \emph{robustly positive}; IST inflation sign is \emph{fragile}, though its wage sign is relative stable.
		\item \textbf{Consumption:} estimated impact is negative; a positive impact needs jointly low $\kappa_p,\psi_\pi,\theta_w$ (a stable real rate).
		\item \textbf{SW (2007):} fits hours, but shrinks the ID shock to a marginal role and flips IST to inflationary.
		\item \textbf{i.i.d.\ ID:} importance falls (output $\approx23\%$, investment $\approx15\%$) but the shock remains a meaningful output driver.
	\end{enumerate}

	\begin{block}{Takeaways}
		\begin{enumerate}
			\item[\textbf{(i)}] SW model shows stronger ability to fit with the data, and delivers posteriors indicating a stickier economy.
			\item[\textbf{(ii)}] Positive consumption impact require strict parameter values, which is not easy to reproduce in the estimation, but the logic of reducing short-run intertemporal substitution effect works.
		\end{enumerate}
	\end{block}
\end{frame}

\begin{frame}[allowframebreaks]{References}
	\bibliography{reference}
\end{frame}

\end{document}
